% Options for packages loaded elsewhere
\PassOptionsToPackage{unicode}{hyperref}
\PassOptionsToPackage{hyphens}{url}
\PassOptionsToPackage{dvipsnames,svgnames,x11names}{xcolor}
%
\documentclass[
  10pt,
]{article}
\usepackage{amsmath,amssymb}
\usepackage{iftex}
\ifPDFTeX
  \usepackage[T1]{fontenc}
  \usepackage[utf8]{inputenc}
  \usepackage{textcomp} % provide euro and other symbols
\else % if luatex or xetex
  \usepackage{unicode-math} % this also loads fontspec
  \defaultfontfeatures{Scale=MatchLowercase}
  \defaultfontfeatures[\rmfamily]{Ligatures=TeX,Scale=1}
\fi
\usepackage{lmodern}
\ifPDFTeX\else
  % xetex/luatex font selection
  \setmainfont[]{DejaVu Serif}
  \setsansfont[]{DejaVu Sans}
  \setmonofont[]{DejaVu Sans Mono}
  \setmathfont[]{Latin Modern Math}
\fi
% Use upquote if available, for straight quotes in verbatim environments
\IfFileExists{upquote.sty}{\usepackage{upquote}}{}
\IfFileExists{microtype.sty}{% use microtype if available
  \usepackage[]{microtype}
  \UseMicrotypeSet[protrusion]{basicmath} % disable protrusion for tt fonts
}{}
\makeatletter
\@ifundefined{KOMAClassName}{% if non-KOMA class
  \IfFileExists{parskip.sty}{%
    \usepackage{parskip}
  }{% else
    \setlength{\parindent}{0pt}
    \setlength{\parskip}{6pt plus 2pt minus 1pt}}
}{% if KOMA class
  \KOMAoptions{parskip=half}}
\makeatother
\usepackage{xcolor}
\usepackage[margin=2.0cm]{geometry}
\usepackage{color}
\usepackage{fancyvrb}
\newcommand{\VerbBar}{|}
\newcommand{\VERB}{\Verb[commandchars=\\\{\}]}
\DefineVerbatimEnvironment{Highlighting}{Verbatim}{commandchars=\\\{\}}
% Add ',fontsize=\small' for more characters per line
\newenvironment{Shaded}{}{}
\newcommand{\AlertTok}[1]{\textcolor[rgb]{1.00,0.00,0.00}{\textbf{#1}}}
\newcommand{\AnnotationTok}[1]{\textcolor[rgb]{0.38,0.63,0.69}{\textbf{\textit{#1}}}}
\newcommand{\AttributeTok}[1]{\textcolor[rgb]{0.49,0.56,0.16}{#1}}
\newcommand{\BaseNTok}[1]{\textcolor[rgb]{0.25,0.63,0.44}{#1}}
\newcommand{\BuiltInTok}[1]{\textcolor[rgb]{0.00,0.50,0.00}{#1}}
\newcommand{\CharTok}[1]{\textcolor[rgb]{0.25,0.44,0.63}{#1}}
\newcommand{\CommentTok}[1]{\textcolor[rgb]{0.38,0.63,0.69}{\textit{#1}}}
\newcommand{\CommentVarTok}[1]{\textcolor[rgb]{0.38,0.63,0.69}{\textbf{\textit{#1}}}}
\newcommand{\ConstantTok}[1]{\textcolor[rgb]{0.53,0.00,0.00}{#1}}
\newcommand{\ControlFlowTok}[1]{\textcolor[rgb]{0.00,0.44,0.13}{\textbf{#1}}}
\newcommand{\DataTypeTok}[1]{\textcolor[rgb]{0.56,0.13,0.00}{#1}}
\newcommand{\DecValTok}[1]{\textcolor[rgb]{0.25,0.63,0.44}{#1}}
\newcommand{\DocumentationTok}[1]{\textcolor[rgb]{0.73,0.13,0.13}{\textit{#1}}}
\newcommand{\ErrorTok}[1]{\textcolor[rgb]{1.00,0.00,0.00}{\textbf{#1}}}
\newcommand{\ExtensionTok}[1]{#1}
\newcommand{\FloatTok}[1]{\textcolor[rgb]{0.25,0.63,0.44}{#1}}
\newcommand{\FunctionTok}[1]{\textcolor[rgb]{0.02,0.16,0.49}{#1}}
\newcommand{\ImportTok}[1]{\textcolor[rgb]{0.00,0.50,0.00}{\textbf{#1}}}
\newcommand{\InformationTok}[1]{\textcolor[rgb]{0.38,0.63,0.69}{\textbf{\textit{#1}}}}
\newcommand{\KeywordTok}[1]{\textcolor[rgb]{0.00,0.44,0.13}{\textbf{#1}}}
\newcommand{\NormalTok}[1]{#1}
\newcommand{\OperatorTok}[1]{\textcolor[rgb]{0.40,0.40,0.40}{#1}}
\newcommand{\OtherTok}[1]{\textcolor[rgb]{0.00,0.44,0.13}{#1}}
\newcommand{\PreprocessorTok}[1]{\textcolor[rgb]{0.74,0.48,0.00}{#1}}
\newcommand{\RegionMarkerTok}[1]{#1}
\newcommand{\SpecialCharTok}[1]{\textcolor[rgb]{0.25,0.44,0.63}{#1}}
\newcommand{\SpecialStringTok}[1]{\textcolor[rgb]{0.73,0.40,0.53}{#1}}
\newcommand{\StringTok}[1]{\textcolor[rgb]{0.25,0.44,0.63}{#1}}
\newcommand{\VariableTok}[1]{\textcolor[rgb]{0.10,0.09,0.49}{#1}}
\newcommand{\VerbatimStringTok}[1]{\textcolor[rgb]{0.25,0.44,0.63}{#1}}
\newcommand{\WarningTok}[1]{\textcolor[rgb]{0.38,0.63,0.69}{\textbf{\textit{#1}}}}
\usepackage{longtable,booktabs,array}
\usepackage{calc} % for calculating minipage widths
% Correct order of tables after \paragraph or \subparagraph
\usepackage{etoolbox}
\makeatletter
\patchcmd\longtable{\par}{\if@noskipsec\mbox{}\fi\par}{}{}
\makeatother
% Allow footnotes in longtable head/foot
\IfFileExists{footnotehyper.sty}{\usepackage{footnotehyper}}{\usepackage{footnote}}
\makesavenoteenv{longtable}
\setlength{\emergencystretch}{3em} % prevent overfull lines
\providecommand{\tightlist}{%
  \setlength{\itemsep}{0pt}\setlength{\parskip}{0pt}}
\setcounter{secnumdepth}{-\maxdimen} % remove section numbering
% doc_header.tex -- LaTeX preamble additions used by build_docs.sh for every
% Markdown -> LaTeX -> PDF conversion (the documents of the folder Revision_Lord_Gravity).
\usepackage{etoolbox}
% longtable may not be loaded by pandoc when a document has no table
\usepackage{longtable}
% bold vectors: \bm -> \symbf under unicode-math (XeLaTeX), bm package otherwise
\makeatletter
\@ifpackageloaded{unicode-math}{\providecommand{\bm}[1]{\symbf{#1}}}{\usepackage{bm}}
\makeatother
% smaller type in long tables so that wide audit tables fit the text width;
% inside a table a line break is allowed after an underscore of plain text
% (pandoc writes it as \_), so that long identifiers wrap inside their column
% rows are set a little apart, so that symbols with indices do not touch
\AtBeginEnvironment{longtable}{\footnotesize\renewcommand{\arraystretch}{1.45}%
  \let\sptunderscore\_\def\_{\sptunderscore\allowbreak}}
\setlength{\LTcapwidth}{\textwidth}
% A short table is kept on one page.  The build (code/scripts/keep_tables.py) ends the rows of such a table with \\*,
% which forbids a page break after a row, and sets \sptkeptpenalty to 10000 in front of the table.  This penalty is
% put here in the place of the breakpoint (penalty 0) that longtable leaves after every table, so that the paragraph
% that follows a kept table, its caption, begins on the same page; it is set back to 0 after the table.
\newcount\sptkeptpenalty
\sptkeptpenalty=0
\makeatletter
\patchcmd{\endlongtable}{\endgraf\penalty\z@\addvspace\LTpost}%
  {\endgraf\penalty\sptkeptpenalty\addvspace\LTpost\global\sptkeptpenalty=\z@}%
  {}{\PackageError{doc_header}{the end of longtable could not be patched}{The tables can still be typeset, but a short table may be separated from its caption.}}
\makeatother
% allow line breaks after slashes and hyphens inside table cells
\usepackage{xurl}
% code blocks: break over-long lines instead of running into the margin.
% Plain ``` blocks are typeset by pandoc as {verbatim}, highlighted ones
% (```bash) as {Highlighting}; both are re-declared as breakable fancyvrb
% environments.  The continuation line is marked with a small hooked arrow.
\usepackage{fvextra}
\RecustomVerbatimEnvironment{verbatim}{Verbatim}{breaklines,breakanywhere,fontsize=\small}
\makeatletter
\@ifundefined{Highlighting}{}{%
  \RecustomVerbatimEnvironment{Highlighting}{Verbatim}{breaklines,breakanywhere,fontsize=\small,commandchars=\\\{\}}}
\makeatother
% inline code (pandoc writes it as \texttt{...}): never hyphenate -- an inserted
% hyphen could be read as part of a file name -- but allow a line break after
% an underscore, slash, comma, colon or equals sign, so that long paths and
% column lists wrap inside the text width and inside table columns.
% Inside a formula (e.g. the name of a notebook symbol in an equation) the same
% command must first leave math mode, which \text of amsmath does.
\ExplSyntaxOn
\tl_new:N \l__spt_code_tl
\cs_new_protected:Npn \__spt_texttt:n #1
  {
    \tl_set:Nn \l__spt_code_tl {#1}
    \regex_replace_all:nnN { ([\/\,\:\=]) } { \1 \c{allowbreak} } \l__spt_code_tl
    \regex_replace_all:nnN { \c{_} } { \c{_} \c{allowbreak} } \l__spt_code_tl
    \group_begin: \ttfamily \hyphenchar \font = -1 \scan_stop: \l__spt_code_tl \group_end:
  }
\RenewDocumentCommand \texttt { m }
  {
    \mode_if_math:TF { \text { \__spt_texttt:n {#1} } } { \__spt_texttt:n {#1} }
  }
\ExplSyntaxOff
% A figure is written in the Markdown source as an image paragraph followed by
% a caption paragraph ("**Figure S1.** ..."), so that the caption is also
% visible on GitHub.  An image that starts a paragraph is centered.
\usepackage{graphicx}% pandoc loads it only for documents that contain an image
\NewCommandCopy\sptincludegraphics\includegraphics
\RenewDocumentCommand\includegraphics{O{} m}{%
  \ifvmode\begin{center}\sptincludegraphics[#1]{#2}\end{center}%
  \else\sptincludegraphics[#1]{#2}\fi}
% keep figures where they are referenced
\usepackage{float}
\floatplacement{figure}{H}
% running header with the document title
\usepackage{fancyhdr}
\pagestyle{fancy}
\fancyhf{}
\fancyhead[L]{\small PROJECT spacetimes, Revision\_Lord\_Gravity}
\fancyhead[R]{\small\thepage}
\renewcommand{\headrulewidth}{0.4pt}
% 16x16 matrices: tighter columns so that a full matrix fits the text width
\setlength{\arraycolsep}{3.2pt}
% The documents of this project contain long tagged symbols, file paths and
% identifiers that cannot be hyphenated.  A paragraph is therefore allowed to
% have looser lines before TeX gives up and lets a line stick out of the margin.
\AtBeginDocument{\tolerance=2000 \emergencystretch=5em \hbadness=2000}
\ifLuaTeX
  \usepackage{selnolig}  % disable illegal ligatures
\fi
\IfFileExists{bookmark.sty}{\usepackage{bookmark}}{\usepackage{hyperref}}
\IfFileExists{xurl.sty}{\usepackage{xurl}}{} % add URL line breaks if available
\urlstyle{same}
\hypersetup{
  pdftitle={The provenance of the notebook lord\_gravity.nb},
  pdfauthor={Prepared for Patrick Lee Nash, Ph.D.},
  colorlinks=true,
  linkcolor={blue},
  filecolor={Maroon},
  citecolor={blue},
  urlcolor={blue},
  pdfcreator={LaTeX via pandoc}}

\title{The provenance of the notebook lord\_gravity.nb}
\usepackage{etoolbox}
\makeatletter
\providecommand{\subtitle}[1]{% add subtitle to \maketitle
  \apptocmd{\@title}{\par {\large #1 \par}}{}{}
}
\makeatother
\subtitle{Provenance document of the notebook, PROJECT spacetimes,
folder Revision\_Lord\_Gravity (task C23, stage 8)}
\author{Prepared for Patrick Lee Nash, Ph.D.}
\date{November 22, 2026}

\begin{document}
\maketitle

\hypertarget{what-this-document-is}{%
\section{1. What this document is}\label{what-this-document-is}}

This document records how the notebook \texttt{lord\_gravity.nb} of the
folder \texttt{Revision\_Lord\_Gravity} was written, built and
evaluated, what it contains and what it leaves out, and how its results
were checked. Every number below is taken by the program
\texttt{code/scripts/write\_lg\_notebook\_fragments.py} from the log
that the evaluation wrote (\texttt{data/hck\_lg\_notebook\_log.txt}) and
from the report of the builder
(\texttt{data/hck\_lg\_notebook\_build.txt}); none is typed by hand.
``I'' is the assistant that did the work.

\hypertarget{what-the-notebook-is-for}{%
\section{2. What the notebook is for}\label{what-the-notebook-is-for}}

The manuscript of the folder, \texttt{lord\_gravity.md}, is checked by
programs in three languages: Rust, Python and the Wolfram Language. The
nine scripts in the Wolfram Language (\texttt{wolfram/source/lg\_*.wls})
decide most of their statements in exact arithmetic: identities between
symbolic matrices, symbolic functions and series, which a floating-point
program can only test at sample points. The plan of the folder (stage 8)
asks for a notebook of the folder in the Wolfram Language, with every
formula and check of the manuscript and its remarks. This version
contains every check that the Wolfram Language makes in the folder and
the remarks; it does not contain every formula of the manuscript (the
next paragraph but one says which it leaves out). It is written for a
student: it runs the code of all nine scripts, in the order of the
manuscript, and puts before each check a text that says, in words and
with the symbols of the manuscript, what the check decides and why it
matters. A student can evaluate it, read it, change an input and watch a
check fail.

The notebook has a title and an introduction (what it is, how to use it,
how to read it, the checks and their counterparts in Rust and Python,
what it does not repeat), a Chapter 0 that defines the register of the
checks, nine chapters (one for each script), a chapter with the remarks
R1 to R16 of the manuscript in brief, and a final table of all checks.

\textbf{What it does not repeat.} The long numerical computations of the
manuscript, the evolutions in time of Parts V and VI, are done by the
Rust programs with the integrator CVODE of SUNDIALS and take from
minutes to forty minutes; the notebook does not repeat them. It repeats
the numbers that a Wolfram script recomputes independently: three
quasi-bound states and three greybody factors of Part IV (checks WS-7
and WS-8), found again in arithmetic of 32 digits by an integration in
the radius with \texttt{NDSolve} (its automatic method, which controls
the error; this is not an integration in time).

\hypertarget{the-files}{%
\section{3. The files}\label{the-files}}

The files of the notebook, each with its SHA-256 checksum:

\begin{itemize}
\tightlist
\item
  \texttt{lord\_gravity.nb}: the notebook, evaluated in the front end of
  Wolfram (inputs, printed checks, the table of all checks);\\
  \texttt{27148f94481551c9290cd32f29b562e5ff61ac74964e732526b4ee731f37c1a2}
\item
  \texttt{lord\_gravity\_notebook.pdf}: the evaluated notebook printed
  by the front end, for readers without Wolfram;\\
  \texttt{3f1935dc945f28be71979f03efc2f68d660d101ad0fb14fbdf9f0d63bbf2498e}
\item
  \texttt{wolfram/notebook/lord\_gravity.texts.wl}: the hand-written
  texts of the notebook, with the marker lines that stand for the code
  of the scripts;\\
  \texttt{0b68045712c74bbd84e13c0425afd943c92b8a19f80901905d86cb633e7c22e0}
\item
  \texttt{wolfram/notebook/lord\_gravity.source.wl}: the complete source
  of the notebook, written by
  \texttt{wolfram/tools/make\_notebook\_source.py} from the texts and
  the nine scripts;\\
  \texttt{93ecb2cad44c9e078a8c11a175a76339052a47c09316fc36fa27d6b309cff7a5}
\item
  \texttt{wolfram/tools/make\_notebook\_source.py}: writes the source
  from the texts and the scripts, and fails if a part of a script is
  left out or used twice;\\
  \texttt{30f10c506349c4b1dfdb33d628aa9e5c84173ac739f5844417267de70699cf88}
\item
  \texttt{wolfram/tools/build\_notebook.wls}: the builder: turns the
  source into the notebook;\\
  \texttt{951a3593a2f5e48590baac1b9007705d0317e8e98f591dbbd2b80b4bd70575ed}
\item
  \texttt{wolfram/tools/lg\_formula.wl}: the reader of formulas: turns
  the TeX of the formulas of the texts into boxes;\\
  \texttt{48d87194b16688989947350b24b19f736977c7644b2dd5da8318895a4d3a9f03}
\item
  \texttt{wolfram/tools/make\_dryrun.py}: writes the code of all input
  cells into one file for a run without the front end;\\
  \texttt{381dee3c1e869d6bbe11097487ce7bd6ef58384d3f91cb7873af1abf546a1925}
\item
  \texttt{wolfram/tools/guirun.wl}: the evaluation of the notebook in
  the front end, in a fresh kernel, with the log;\\
  \texttt{195f4bdc73d58df6546c9f18e16ef8ab08ced5e6891c09f62df075717a5bca4a}
\item
  \texttt{wolfram/tools/gui\_start.wl}: the first of the three calls of
  the evaluation;\\
  \texttt{004868abf8533468f07dd013cdd56ca2bc65efce4752904b194a8c5b9e2bb4e6}
\item
  \texttt{wolfram/tools/gui\_go.wl}: the second call;\\
  \texttt{1ffdf2aea766e923a10abe53fb992d64e05b6328b77dfe1a96050cb7800bb2ad}
\item
  \texttt{wolfram/tools/gui\_status.wl}: the third call;\\
  \texttt{6529fcec46f622bcf3098ac60ee9ea4e95d6abaa007a18a7bda9a1ffc7d34819}
\end{itemize}

The notebook also contains, in closed cells (cells that show only a thin
bracket and are left out of the printed copy), the three data files that
the checks compare with: \texttt{data/lg\_schw\_qbs.csv} and
\texttt{data/lg\_schw\_grey.csv} (the quasi-bound states and the
greybody factors of the Rust program \texttt{lg\_schw}),
\texttt{data/lg\_pair\_levels.csv} (the exact levels of the pair,
written by \texttt{python/lg\_pair.py}), and the fragment file
\texttt{provenance/generated/lg\_examples.md} (the printed numbers of
the appendices); each closed cell names its file and the SHA-256
checksum that the file had when the source was written.

\hypertarget{how-the-notebook-is-made}{%
\section{4. How the notebook is made}\label{how-the-notebook-is-made}}

\textbf{From the scripts to the source.} The texts of the notebook are
written by hand in \texttt{wolfram/notebook/lord\_gravity.texts.wl}, in
the format of the builder: every cell begins with a marker line such as
\texttt{(*\ ::Text::\ *)} or \texttt{(*\ ::Input::\ *)}. Where the code
of a script belongs, the texts have a marker line instead, for example
one that stands for the section of the script \texttt{lg\_lord.wls} that
makes the check WL-1. The program
\texttt{wolfram/tools/make\_notebook\_source.py} replaces every marker
line by an input cell with that code, taken from the script, and writes
the complete source \texttt{wolfram/notebook/lord\_gravity.source.wl}.
It fails if a section of a script is left out or used twice, so the
notebook runs all of the code of the nine scripts, and only that code.
It leaves out the lines of each script that read its arguments, count
its checks and write its report (the notebook has one register of the
checks for all chapters), and it changes the code of three scripts so
that the notebook reads no file: \texttt{lg\_schw.wls},
\texttt{lg\_pair.wls} and \texttt{lg\_worked.wls} read the text of their
data file from a closed cell instead of the file. Each chapter runs in a
context of its own (\texttt{LGNB\textasciigrave{}Lord\textasciigrave{}},
\texttt{LGNB\textasciigrave{}Curved\textasciigrave{}}, and so on), so
that the symbols of one script cannot change those of another, and the
assumptions that one script gives to \texttt{Simplify} are reset at the
end of its chapter.

\textbf{The register of the checks.} Chapter 0 defines the functions
\texttt{check} and \texttt{checkDev} once, for all chapters. They print
the same line as the scripts print into their reports, and they append
every check to a list that the last chapter prints as a table and that
the evaluation writes into its log.

\textbf{From the source to the notebook.} The builder
\texttt{wolfram/tools/build\_notebook.wls} turns the source into a
notebook: every input cell is parsed by the front end of Wolfram (so
that comments, line breaks and indentation stay as written), and every
formula of a text, written in TeX between dollar signs, is turned into a
two-dimensional formula by the reader of formulas
\texttt{wolfram/tools/lg\_formula.wl}, which reports every command that
it cannot read. The builder, the reader of formulas and the programs of
the evaluation were written for earlier notebooks of the project, the
last of them that of the folder \texttt{Revision\_SM\_Higgs}, and are
copied here with the names of their contexts and of their work folder
changed.

\textbf{The dry run.} Before the notebook is built,
\texttt{wolfram/tools/make\_dryrun.py} writes the code of all input
cells into one file, which \texttt{wolframscript} evaluates in a plain
kernel, without the front end. This finds errors of the code in seconds.

\textbf{The evaluation.} The notebook is evaluated in the front end of
the running Wolfram application, exactly as by the menu command
Evaluation \textgreater{} Evaluate Notebook, but in a fresh kernel, in
three calls: \texttt{gui\_start.wl} opens the built notebook and starts
the kernel; \texttt{gui\_go.wl} evaluates all cells and then writes the
register of the checks to a file; \texttt{gui\_status.wl} saves the
evaluated notebook, prints it to PDF, and writes the log (the version,
the time, the counts of the cells, every message and every printed line,
and the register of the checks). On the author's computer these calls,
the dry run and the builder were run by a small program inside the
Wolfram application that executes requests left as files in a work
folder of the author's copy of the repository (requests 351 to 357 of 22
November 2026).

\hypertarget{the-evaluation-of-this-version}{%
\section{5. The evaluation of this
version}\label{the-evaluation-of-this-version}}

The notebook was built with 228 cells (4 of them closed cells of data)
and 422 formulas in its texts, with 0 warnings of the reader of
formulas. It was evaluated in the front end of 15.0.1 for Mac OS X ARM
(64-bit) (July 2, 2026) in 48 seconds: 84 input cells, 84 of them
evaluated, 1 output cell (the table of all checks), 53 printed lines,
and \textbf{0 messages} (no error and no warning). Of its 46 checks, 46
passed.

\textbf{Table 1.} The chapters of the notebook, the checks that each
makes (with the number of the input line that printed the check), and
their results.

\sptkeptpenalty=10000
\begin{longtable}[]{@{}
  >{\raggedright\arraybackslash}p{(\columnwidth - 8\tabcolsep) * \real{0.0727}}
  >{\raggedright\arraybackslash}p{(\columnwidth - 8\tabcolsep) * \real{0.2909}}
  >{\raggedright\arraybackslash}p{(\columnwidth - 8\tabcolsep) * \real{0.1636}}
  >{\raggedright\arraybackslash}p{(\columnwidth - 8\tabcolsep) * \real{0.3455}}
  >{\raggedright\arraybackslash}p{(\columnwidth - 8\tabcolsep) * \real{0.1273}}@{}}
\toprule\noalign{}
\begin{minipage}[b]{\linewidth}\raggedright
chapter
\end{minipage} & \begin{minipage}[b]{\linewidth}\raggedright
part of the manuscript
\end{minipage} & \begin{minipage}[b]{\linewidth}\raggedright
script
\end{minipage} & \begin{minipage}[b]{\linewidth}\raggedright
checks (In line of the notebook)
\end{minipage} & \begin{minipage}[b]{\linewidth}\raggedright
result
\end{minipage} \\
\midrule\noalign{}
\endhead
\bottomrule\noalign{}
\endlastfoot
1 & Parts I and II: Lord's paper and the association &
\texttt{lg\_lord.wls} & WL-1 (19), WL-2 (31), WL-3 (37), WL-4 (42), WL-5
(47), WL-6 (59), WL-7 (67) & 7 of 7 pass \\*
2 & Part III: the brochon in a curved space-time &
\texttt{lg\_curved.wls} & WG-1 (97), WG-2 (113), WG-3 (126), WG-4 (136),
WG-5 (152), WG-6 (187) & 6 of 6 pass \\*
3 & Part IV: the Schwarzschild field & \texttt{lg\_schw.wls} & WS-1
(222), WS-2 (230), WS-3 (237), WS-4 (249), WS-5 (256), WS-9 (260), WS-6
(264), WS-7 (282), WS-8 (292) & 9 of 9 pass \\*
4 & Part V: the Vaidya field & \texttt{lg\_vaidya.wls} & WV-1 (315),
WV-2 (326), WV-3 (329), WV-4 (340), WV-4b (340), WV-5 (353) & 6 of 6
pass \\*
5 & Part VI, Section 29: the pair & \texttt{lg\_pair.wls} & WP-1 (362),
WP-2 (370) & 2 of 2 pass \\*
6 & Part VI, Section 30: the charged black hole &
\texttt{lg\_charge.wls} & WC-1 (382), WC-2 (387), WC-3 (395) & 3 of 3
pass \\*
7 & Part VI, Section 32: the fall of the pair & \texttt{lg\_tidal.wls} &
WT-1 (407), WT-2 (411), WT-3 (430), WT-4 (438), WT-5 (458) & 5 of 5
pass \\*
8 & Part VI, Section 33: annihilation & \texttt{lg\_annih.wls} & WA-1
(486), WA-2 (488), WA-3 (500) & 3 of 3 pass \\*
9 & the worked examples & \texttt{lg\_worked.wls} & WW-1 (529), WW-2
(542), WW-3 (547), WW-4 (563), WW-5 (578) & 5 of 5 pass \\
\end{longtable}

Because the notebook runs the code of the scripts, it must reach their
results. The program
\texttt{code/scripts/write\_lg\_notebook\_fragments.py} compares the
line that each check printed in the notebook with the line of the same
check in the report of its script
(\texttt{data/hck\_lg\_*\_checks\_wolfram.txt}): 46 of the 46 lines are
equal, character for character, including the measured deviations of the
numerical checks WS-7 and WS-8. Its complete report
(\texttt{data/hck\_lg\_notebook\_checks.txt}):

\begin{verbatim}
PROJECT spacetimes, Revision_Lord_Gravity: the checks of the provenance of the notebook lord_gravity.nb (stage 8 of task C23; code/scripts/write_lg_notebook_fragments.py)

PASS LN-1  the notebook was evaluated in the front end of 15.0.1 for Mac OS X ARM (64-bit) (July 2, 2026) in 48 s: 84 input cells, 84 of them evaluated; 0 messages; 46 checks registered, 46 passed
PASS LN-2  the 46 checks of the nine scripts were each made once in the notebook, and no other check
PASS LN-3  the line that each check printed in the notebook is character for character the line of the same check in the report of its script: 46 of 46
PASS LN-4  the builder made 228 cells (4 of them closed cells of data) with 422 formulas in the text cells, and reported 0 warnings of the reader of formulas

VERDICT: ALL PASS (4 of 4)
\end{verbatim}

\hypertarget{corrections-made-on-the-way}{%
\section{6. Corrections made on the
way}\label{corrections-made-on-the-way}}

\begin{itemize}
\tightlist
\item
  \textbf{A message of the script \texttt{lg\_tidal.wls}.} The first dry
  run (request 351) passed all checks but printed one message,
  ClebschGordan::tri, in the check WT-1: the 3-j symbol of the angular
  momenta 0, 2, 0 violates the triangle condition, and
  \texttt{ThreeJSymbol} returns 0 for it, which is the correct coupling,
  but with a message. The script had printed the same message when it
  was first run (request 349); its report does not show messages. The
  script now sets such a 3-j symbol to 0 without calling
  \texttt{ThreeJSymbol}; it was run again (request 352) and its report
  is unchanged, byte for byte. The notebook printed no message.
\item
  \textbf{A formula that the reader could not read.} The first build
  (request 353) warned that the reader of formulas does not know the
  command for the box of the wave operator, in the text of the check
  WG-4; the text also wrote the Dirac operator as a slashed letter,
  which the reader of formulas does not set as one symbol. The text now
  writes the Dirac operator and its square out with the covariant
  derivative and the inverse metric. The second build (request 354) had
  no warning.
\item
  \textbf{The notation of the texts.} The texts of the first build used,
  in Parts III to VI, short letters of the scripts (for example the
  letters of the speed of the falling observers and of the mass of the
  black hole without their suffixes). For the second build they were
  rewritten with the symbols of the manuscript, and the names of the
  code are given beside them.
\end{itemize}

\hypertarget{how-to-reproduce}{%
\section{7. How to reproduce}\label{how-to-reproduce}}

From the folder \texttt{Revision\_Lord\_Gravity}, with Python 3 and
Wolfram 15:

\begin{Shaded}
\begin{Highlighting}[]
\ExtensionTok{python3}\NormalTok{ wolfram/tools/make\_notebook\_source.py                    }\CommentTok{\# the source, from the texts and the nine scripts}
\ExtensionTok{python3}\NormalTok{ wolfram/tools/make\_dryrun.py dry.wls wolfram/notebook/lord\_gravity.source.wl}
\ExtensionTok{wolframscript} \AttributeTok{{-}file}\NormalTok{ dry.wls                                     }\CommentTok{\# the dry run, in a plain kernel (about 40 s)}
\ExtensionTok{wolframscript} \AttributeTok{{-}file}\NormalTok{ wolfram/tools/build\_notebook.wls wolfram/notebook/lord\_gravity.source.wl lord\_gravity.nb }\DataTypeTok{\textbackslash{}}
\NormalTok{    wolfram/tools/lg\_formula.wl                                 }\CommentTok{\# the notebook (not evaluated)}
\end{Highlighting}
\end{Shaded}

Then open \texttt{lord\_gravity.nb} in Wolfram and choose Evaluation
\textgreater{} Evaluate Notebook; or evaluate it in a fresh kernel with
the three calls of Section 4, after changing the paths of the work
folder in \texttt{wolfram/tools/guirun.wl} and
\texttt{wolfram/tools/gui\_start.wl}. Finally
\texttt{python3\ code/scripts/write\_lg\_notebook\_fragments.py} reads
the log and writes the fragments and the checks of this document.

\end{document}
